\subsection{UPPER CENTRAL SERIES}

Let \(\mathfrak{g}\) be a Lie algebra and \(\mathbf{P}\) a subset of \(\mathfrak{g}\) . The centralizer of \(\mathbf{P}\) in \(\mathfrak{g}\) is the set of elements of \(\mathfrak{g}\) which are permutable with those of \(\mathbf{P}\) . This centralizer is the intersection of the kernels of the ad \(y\) , where \(y\) runs through \(\mathbf{P}\) ; it is therefore a subalgebra of \(\mathfrak{g}\) .

PROPOSITION 5. Let g be a Lie algebra and a an ideal (resp. a characteristic ideal) of g. The centralizer \(\mathfrak{a}'\) of a in g is an ideal (resp. a characteristic ideal) of g.

Let \(D\) be an inner derivation (resp. a derivation) of \(g\) . If \(x \in a'\) and \(y \in a\) , then

\[
[ \mathrm{D} x, y ] = \mathrm{D} ([ x, y ]) - [ x, \mathrm{D} y ] = 0;
\]

hence \(\mathrm{D}x\in \mathfrak{a}^{\prime}\) . Hence the proposition.

Let g be a Lie algebra. The centralizer of g in g is called the centre of g, that is the characteristic ideal of \(x \in g\) such that \([x, y] = 0\) for all \(y \in g\) . The centre of g is the kernel of the homomorphism \(x \mapsto ad x\) .

The upper central series of \(\mathfrak{g}\) is the increasing sequence \(\mathcal{C}_{0}\mathfrak{g},\mathcal{C}_{1}\mathfrak{g},\ldots\) of characteristic ideals of \(\mathfrak{g}\) defined inductively as follows: (1) \(\mathcal{C}_{0}\mathfrak{g} = \{0\}\) ; (2) \(\mathcal{C}_{p + 1}\mathfrak{g}\) is the inverse image under the canonical mapping of \(\mathfrak{g}\) onto \(\mathfrak{g} / \mathcal{C}_{p}\mathfrak{g}\) of the centre of \(\mathfrak{g} / \mathcal{C}_{p}\mathfrak{g}\) .

The ideal \(\mathcal{C}_{1}\mathfrak{g}\) is the centre of \(\mathfrak{g}\) .

